\begin{frame}{The peel is a vector over levels --- and one offset relocates instead of peeling}

The retime matrix is indexed by \key{(op-class, level)}, so \emph{any} level can carry an offset: an
outer tile level, the reduction chunk, the substep. \texttt{PeelDepths} therefore reports
$M_\ell$ per level, the per-op-class offsets, the substep request, the derived reach, and every
negative offset separately as a deferral.

\begin{itemize}
  \item \texttt{PeelDepths.M} is deliberately \key{narrow} --- the chunk level, the only one whose prologue/drain the emitter can build --- and sits beside \texttt{peeled\_levels()}, which reports \emph{all} of them. That pairing turns a silent drop into a loud refusal.
  \item It used to return an \texttt{int} keyed on the chunk. The tree stayed well formed, the ledger came back empty, and the certificate \bad{certified a schedule nobody requested}.
\end{itemize}

\textbf{The cross-level boundary term.} An offset at outer level $\ell$ whose op-class lives strictly
\emph{inside} $\ell$ does not peel $\ell$ --- it relocates the leading $\delta$ instances into the
previous $\ell$-trip's drain. Three conditions decide it, and the third separates these two:

\begin{center}\scriptsize
\begin{tabular}{lll}
\toprule
offset & home level & verdict \\
\midrule
\texttt{off(copy, gtile)} & \texttt{iter} --- an outer, looped level & \good{relocation}: whole iterations move \\
\texttt{off(copy, iter)}  & \texttt{M\_split} --- an inner unrolled axis & \bad{ordinary prefetch}: the fan stays put \\
\bottomrule
\end{tabular}
\end{center}

\tsub{Without condition 3 the two are conflated --- the first implementation reported a boundary term for every shipping region-split kernel, and the emitter refuses on a non-empty hoist list, so it refused them all.}
\end{frame}
